\subsection{BIGEBRAS AND SKEW-BIGEBRAS}

DEFINITION 3. A graded bigebra (resp. skew graded bigebra) over a ring A is a set E with a graded A-algebra structure of type N and a graded A-cogebra structure of type N, with the same underlying graded A-module structure and such that:

\begin{enumerate}
\def\labelenumi{(\arabic{enumi})}
\item
  The A-algebra E is associative and unital.
\item
  The A-cogebra E is coassociative and counital.
\item
  The coproduct \(c: \mathbf{E} \to \mathbf{E} \otimes_{\mathbf{A}} \mathbf{E}\) is a homomorphism of the graded algebra \(\mathbf{E}\) into the graded algebra \(\mathbf{E} \otimes_{\mathbf{A}} \mathbf{E}\) (resp. graded algebra \(\mathbf{E}^{\mathfrak{g}} \otimes_{\mathbf{A}} \mathbf{E}\) (cf.~§ 4, no. 7)).
\item
  The counit \(\gamma\) of E is a homomorphism of the graded algebra E into the algebra A (with the trivial graduation) such that, if e denotes the unit element of the A-algebra \(\mathbf{E}, \gamma(e) = 1\) .
\end{enumerate}

If E is a graded bigebra whose graduation is trivial, E is called simply a bigebra. A graded bigebra is called commutative (resp. cocommutative) if the underlying algebra is commutative (resp. if the underlying cogebra is cocommutative); a skew graded bigebra is called anticommutative (resp. anticocommutative) if the underlying graded algebra is anticommutative (resp. if the underlying graded cogebra is anticocommutative).

It follows from Definition 3 and no. 2, Proposition 5 that, for a graded bigebra or a skew graded bigebra E,

\[
\left\{ \begin{array}{c} c (e) = e \otimes e \\ c (x) \equiv x \otimes e + e \otimes x (\mathrm{mod.} \mathrm{E} _ {+} \otimes \mathrm{E} _ {+}) \quad \text { for } x \in \mathrm{E} _ {+} = \bigoplus_ {n \geqslant 1} \mathrm{E} _ {n}. \end{array} \right.\tag{22}
\]

If E and F are two graded bigebras (resp. two skew graded bigebras), a mapping \(\phi: \mathrm{E} \to \mathrm{F}\) is called a graded bigebra morphism (resp. skew graded bigebra morphism) if: (1) \(\phi\) is a graded algebra morphism (and hence maps the unit element of E to the unit element of F); (2) \(\phi\) is a graded cogebra morphism such that, if \(\gamma\) and \(\gamma'\) are the respective counits of E and F, then \(\gamma = \gamma' \circ \phi\) .

Examples. (1) Let S be a monoid with identity element u, so that the algebra \(E = A^{(S)}\) of the monoid S over A admits the unit element \(e_{u}\) (§2, no. 6); it has been seen on the other hand that E has canonically a coassociative cocommutative A-cogebra structure with a counit \(\gamma\) such that \(\gamma(e_{s}) = 1\) for all \(s \in S\) (no. 1, Example 4 and no. 2, Examples 4, 5 and 6). The formula \(c(e_{s}) = e_{s} \otimes e_{s}\) giving the coproduct shows also immediately that c is an algebra homomorphism. Thus a cocommutative bigebra structure has been defined on E and E, with this structure, is called the bigebra of the monoid S over A.

If T is another monoid with unit element \(v, f: \mathrm{S} \to \mathrm{T}\) a homomorphism such that \(f(u) = v\) and \(f_{(\mathbf{A})}: \mathrm{A}^{(\mathrm{S})} \to \mathrm{A}^{(\mathrm{T})}\) the A-algebra homomorphism derived from \(f(\S 2, \text{no. } 6)\) , it is immediately verified that \(f_{(\mathbf{A})}\) is a bigebra homomorphism.

\begin{enumerate}
\def\labelenumi{(\arabic{enumi})}
\setcounter{enumi}{1}
\item
  Let M be an A-module. The graded A-algebra (§ 6, no. 1) and graded A-cogebra (no. 1, Example 6) structures defined on S(M) define on this set a commutative cocommutative graded bigebra structure; for we have seen (no. 1, Example 6) that the coproduct on S(M) is an algebra homomorphism and it follows from the definition of the counit \(\gamma\) (no. 2, Example 6) that \(\gamma(1) = 1\) and that \(\gamma\) is an algebra homomorphism of \(E\) into \(A\) .
\item
  Let M be an A-module. We see as in Example 2 that the graded A-algebra (§ 7, no. 1) and graded A-cogebra (no. 1, Example 7) structures on \(\bigwedge\) (M) define on this set an anticommutative anticocommutative skew graded bigebra structure.
\end{enumerate}

Remark. If M is an A-module such that \(M \otimes_{A} M \neq \{0\}\) , the graded A-algebra (§ 5, no. 1) and graded A-cogebra (no. 1, Example 5) structures on T(M) do not define a bigebra structure, for in general

\[
c (x _ {1} x _ {2} y _ {1} y _ {2}) \neq c (x _ {1} x _ {2}) c (y _ {1} y _ {2})
\]

for four elements \(x_{1}, x_{2}, y_{1}, y_{2}\) of M, as formula (3) of no. 1, shows.

